\chapter{Introduction}
\label{chap:intro}

\begingroup
\let\clearpage\relax

Speech enhancement is a critical domain within signal processing, dedicated to refining speech signals captured under adverse conditions marked by noise or degradation. The core objective is to enhance the quality and intelligibility of these speech signals. The Real-time Targeted Speech Enhancement project aims to address the challenge of strengthening speech signals in a real-time concerning clear speech provided by the user.

The overarching aim of speech enhancement is to elevate the clarity and understandability of speech, particularly in challenging environments. In many applications, a fundamental requirement for a speech enhancement system is its ability to operate in real-time with minimal latency, preferably on cost-effective hardware integrated into communication devices. The quality of speech in virtual interactions significantly impacts user experience and effective communication.

This document presents an in-depth exploration of the problem definition, background information, related work, proposed models, the application and future work.

\hyperref[chap:problem-def]{Chapter 2} focuses on the problem definition, outlining the core objectives of the project and the challenges associated with speech enhancement in the context of online communication. The chapter also highlights the importance of leveraging deep learning methods specifically tailored for real-time speech enhancement.

\hyperref[chap:background]{Chapter 3} delves into the fundamental concepts of audio representation, including waveforms, time-frequency representations, and output representations. The chapter also explores the application of deep learning methodologies, such as Convolutional Neural Networks and Recurrent Neural Networks, in speech enhancement and audio processing.

\hyperref[chap:datasets]{Chapter 4} discusses the datasets used in the project, including the Valentini dataset, the DNS datasets, and the LibriSpeech dataset and our own collected dataset. The chapter also provides an overview of the data augmentation techniques used to increase the size of the training data.

\hyperref[chap:metrics]{Chapter 5} presents the evaluation metrics used to assess the performance of the speech enhancement models including both objective and subjective metrics. It introduces the Perceptual Evaluation of Speech Quality (PESQ), the Short-Time Objective Intelligibility (STOI) metrics, Deep Noise Suppression Mean Opinion Score (DNSMOS) and other metrics used in the evaluation of speech enhancement models.

\hyperref[chap:related-work]{Chapter 6} offers a comprehensive review of the latest advancements in speech enhancement, focusing on the state of the art and identifying gaps in the current research. This chapter also explores various speaker embedders and their integration into cutting-edge architectures. Furthermore, it emphasizes the unique contributions and potential impact of the Real-time Targeted Speech Enhancement project on the speech enhancement domain.

\hyperref[chap:models]{Chapter 7} focuses on the proposed models developed as part of the project, including the Personalized Fast FullSubNet (PFFSN) and the Personalized Denoiser (PDenoiser). The chapter provides an in-depth analysis of the models' architectures, training procedures, and evaluation results. It also discusses the performance of the models on the collected dataset and the DNS Challenge 2023.


\hyperref[chap:application]{Chapter 8} provides an overview of the application developed as part of the project, detailing its primary components and functionality. The chapter also discusses the workflow of the application, divided into two main stages: the GUI app and the streamer. The GUI app serves as the interface for user interaction, capturing user inputs and passing them to the subsequent stages of the application. The streamer handles real-time audio input, utilizing a Virtual Audio Cable (VB cable) to capture the user's voice. The chapter also provides a detailed explanation of the application's streamer and features.

\hyperref[chap:future-work]{Chapter 9} outlines potential future directions for the project, including the integration of additional features, optimization of the existing models, and the exploration of new datasets. The chapter also discusses the potential for expanding the application's functionality and enhancing its real-time capabilities.

\endgroup
